\chapter[1985 --- Brown et al.]{1985: Michael Stuart Brown, Joseph L. Goldstein}
\label{ch:1985_Brown}

\section*{Main Contribution}

\textit{Discovered the LDL receptor and its role in cholesterol metabolism, explaining the genetic basis of familial hypercholesterolemia and leading to statin drugs.}

\section{Official Citation}

for their discoveries concerning the regulation of cholesterol metabolism

The 1985 Nobel Prize in Physiology or Medicine was awarded jointly to Michael Stuart Brown and Joseph L. Goldstein for their discoveries concerning the regulation of cholesterol metabolism. Working together at the University of Texas Southwestern Medical Center in Dallas, Brown and Goldstein elucidated the molecular mechanisms by which cells take up and regulate cholesterol, discovered the low-density lipoprotein (LDL) receptor, and explained the genetic basis of familial hypercholesterolemia---a common and potentially fatal genetic disorder. Their work provided the scientific foundation for the development of statin drugs, which have become the most widely prescribed cholesterol-lowering medications in the world and have saved millions of lives by reducing the risk of heart attacks and strokes.

\section{Background and Historical Context}

Cholesterol is a waxy, fat-soluble molecule that is essential for life. It is a critical component of cell membranes, where it helps maintain membrane fluidity and integrity. It is also the precursor of steroid hormones (including cortisol, aldosterone, estrogen, and testosterone), bile acids, and vitamin D. However, when cholesterol levels in the blood become too high, cholesterol can accumulate in the walls of arteries, forming atherosclerotic plaques that narrow and harden the blood vessels. This process---atherosclerosis---is the leading cause of heart attacks and strokes, and it accounts for more deaths worldwide than any other disease.

The relationship between blood cholesterol levels and heart disease had been recognized since the 1950s, when the Framingham Heart Study and other epidemiological studies demonstrated a strong, positive correlation between serum cholesterol levels and the risk of coronary heart disease. However, the molecular mechanisms by which cholesterol levels are regulated in the body were poorly understood. It was not known how cells take up cholesterol from the blood, how cholesterol levels are sensed and adjusted, or what genetic factors determine individual variation in cholesterol levels and heart disease risk.

Michael Stuart Brown was born in Brooklyn, New York, in 1941, and trained in medicine at the University of Pennsylvania. Joseph Lawrence Goldstein was born in Sumter, South Carolina, in 1940, and also trained in medicine at the University of Pennsylvania. Both men completed their medical training at the Massachusetts General Hospital before joining the faculty at the University of Texas Southwestern Medical Center in Dallas in the early 1970s. Their collaboration, which began in 1972, would become one of the most productive and influential partnerships in the history of biomedical research.

Brown and Goldstein's interest in cholesterol metabolism was sparked by their study of patients with familial hypercholesterolemia (FH), a genetic disorder that causes extremely high levels of cholesterol in the blood. Patients with the homozygous form of FH (in which both copies of the defective gene are mutated) have cholesterol levels that are three to six times higher than normal, and they develop severe atherosclerosis and coronary heart disease at a very young age---often in childhood or adolescence. Heterozygous carriers of the FH mutation (who have one normal and one defective copy of the gene) have moderately elevated cholesterol levels and are at increased risk of heart disease in adulthood.

\section{The Discovery/Contribution}

\subsection{The LDL Receptor and the Mechanism of Cholesterol Uptake}

Brown and Goldstein's first major breakthrough came in 1973, when they discovered that cells from patients with FH were defective in their ability to take up low-density lipoprotein (LDL)---the major cholesterol-carrying particle in the blood. Using cultured human fibroblasts (skin cells), they showed that normal cells had specific receptor sites on their surface that bound LDL with high affinity and internalized it by a process called receptor-mediated endocytosis. In patients with FH, these receptors were absent or non-functional, and the cells were unable to take up LDL from the culture medium.

Brown and Goldstein's discovery of the LDL receptor was a landmark in cell biology and biochemistry. They showed that the LDL receptor is a transmembrane glycoprotein that binds the apolipoprotein B-100 component of LDL particles and internalizes them in clathrin-coated vesicles. Once inside the cell, the LDL particles are delivered to lysosomes, where the cholesterol esters are hydrolyzed to free cholesterol. This free cholesterol is then used by the cell for membrane synthesis and other metabolic processes.

The discovery of the LDL receptor explained the molecular basis of familial hypercholesterolemia: patients with FH have mutations in the gene encoding the LDL receptor (located on chromosome 19), which impair the receptor's ability to bind LDL, to be internalized, or to reach the cell surface. The severity of the disease depends on the nature and location of the mutation: some mutations completely abolish receptor function (producing the severe homozygous form of FH), while others reduce receptor activity by 50\% or more (producing the milder heterozygous form).

\subsection{Feedback Regulation of Cholesterol Synthesis}

Brown and Goldstein's second major breakthrough was the elucidation of the feedback mechanism by which cells regulate their cholesterol levels. They showed that when the intracellular concentration of free cholesterol increases---either as a result of LDL uptake or from de novo cholesterol synthesis---the cell responds by: (1) suppressing the synthesis of new cholesterol by inhibiting the enzyme HMG-CoA reductase (the rate-limiting enzyme in the cholesterol biosynthetic pathway); (2) suppressing the synthesis of new LDL receptors, thereby reducing the uptake of LDL from the blood; and (3) activating an enzyme (acyl-CoA:cholesterol acyltransferase, or ACAT) that esterifies free cholesterol for storage within the cell.

This feedback regulation ensures that cells maintain a constant level of intracellular cholesterol, regardless of the amount of cholesterol available in the blood. In normal individuals, this system works efficiently to keep blood cholesterol levels within a narrow, healthy range. However, in patients with FH, the system is disrupted because the defective LDL receptors cannot take up LDL from the blood. As a result, the cells behave as if they are cholesterol-depleted, even when blood cholesterol levels are very high: they continue to synthesize cholesterol and continue to express LDL receptors (which are non-functional), leading to a vicious cycle of ever-increasing blood cholesterol levels.

\subsection{The Sterol Regulatory Element-Binding Proteins}

In the 1990s, Brown and Goldstein made another major discovery: they identified the transcription factors that mediate the feedback regulation of cholesterol synthesis and LDL receptor expression. These transcription factors, called sterol regulatory element-binding proteins (SREBPs), are membrane-bound proteins that are activated when intracellular cholesterol levels fall. When activated, SREBPs are cleaved and released from the membrane, and they migrate to the nucleus, where they activate the transcription of genes involved in cholesterol synthesis (including HMG-CoA reductase) and LDL uptake (including the LDL receptor gene).

The discovery of SREBPs completed the molecular picture of cholesterol regulation that Brown and Goldstein had been building for more than two decades. It revealed a elegant homeostatic system in which cells sense their cholesterol levels, adjust their synthesis and uptake accordingly, and maintain cholesterol homeostasis through a delicate balance of feedback mechanisms.

\section{Impact on Modern Medicine}

The discoveries of Brown and Goldstein have had a significant impact on the treatment of cardiovascular disease---the leading cause of death worldwide. a major practical consequence of their work was the development of statin drugs, which lower blood cholesterol levels by inhibiting HMG-CoA reductase, the enzyme that catalyzes the rate-limiting step in cholesterol biosynthesis.

The first statin, lovastatin, was approved by the U.S. Food and Drug Administration (FDA) in 1987, and it was followed by a succession of more potent statins, including simvastatin, pravastatin, atorvastatin (Lipitor), and rosuvastatin (Crestor). Statins have been shown in large-scale clinical trials to reduce blood LDL cholesterol levels by 30--50\%, to reduce the risk of heart attacks by 25--35\%, and to reduce the risk of strokes by approximately 20\%. Statins are now the most widely prescribed drugs in the world, with an estimated 200 million people taking them worldwide, and they have been credited with preventing millions of heart attacks and strokes.

Brown and Goldstein's work also led to the development of other cholesterol-lowering drugs, including ezetimibe (which inhibits the absorption of cholesterol from the intestine), PCSK9 inhibitors (monoclonal antibodies that increase the number of LDL receptors on the cell surface), and bile acid sequestrants. PCSK9 inhibitors, in particular, represent a direct extension of Brown and Goldstein's work: PCSK9 is a protein that promotes the degradation of the LDL receptor, and by blocking PCSK9, these drugs increase the number of functional LDL receptors on the cell surface, thereby lowering blood LDL cholesterol levels to very low levels.

The genetic studies of Brown and Goldstein also provided the foundation for the field of pharmacogenomics---the study of how genetic variation influences drug response. The identification of the LDL receptor gene and other genes involved in cholesterol metabolism has enabled the development of genetic tests for familial hypercholesterolemia and other lipid disorders, and these tests are now used to identify individuals at high risk of cardiovascular disease who may benefit from early intervention with statins or other cholesterol-lowering therapies.

In addition to their contributions to drug development, Brown and Goldstein's work has had a major impact on public health policy. Their discoveries provided the scientific basis for dietary guidelines recommending the reduction of saturated fat and cholesterol intake, and for the widespread screening of blood cholesterol levels in the general population. The recognition that high LDL cholesterol is a major, modifiable risk factor for heart disease has led to a dramatic reduction in the incidence of cardiovascular disease in many countries over the past three decades.

\section{Legacy}

Michael Brown and Joseph Goldstein continue to work at the University of Texas Southwestern Medical Center, where they have trained generations of scientists and physicians. Their laboratory has been one of the most productive in the history of biomedical research, producing numerous important discoveries beyond the LDL receptor and cholesterol regulation, including the elucidation of the mechanisms of iron metabolism and the discovery of the gene responsible for hemochromatosis.

Brown and Goldstein have also been influential advocates for basic research, arguing that a major medical advances come from investigations motivated by curiosity rather than by immediate practical goals. Their own career is a powerful illustration of this principle: the discovery of the LDL receptor---which led to the development of statins and the prevention of millions of heart attacks---was motivated not by a desire to develop a drug, but by a desire to understand the molecular basis of a genetic disease.

The work of Brown and Goldstein demonstrates the power of combining clinical observation with basic science. By studying a rare genetic disease (familial hypercholesterolemia), they discovered fundamental principles of cell biology that have had a significant impact on the health of millions of people. Their legacy is a legacy of scientific rigor, clinical relevance, and life-saving discovery.


\section{Modern Developments}

Brown and Goldstein's cholesterol work has led to statin drugs, which have prevented millions of heart attacks and strokes. PCSK9 inhibitors provide additional cholesterol lowering. Inclisiran offers twice-yearly dosing. Understanding of LDL receptor biology has led to treatments for familial hypercholesterolemia. Bempedoic acid offers an alternative for statin-intolerant patients. Understanding of lipid metabolism has informed cardiovascular risk prediction and prevention.


\section{Bibliography}

\begin{thebibliography}{9}
\bibitem{nobel-1985}
Nobel Prize Outreach, ``The Nobel Prize in Physiology or Medicine 1985,''\emph{NobelPrize.org}.\\
\url{https://www.nobelprize.org/prizes/medicine/1985/summary/}

\bibitem{lecture-1985}
Michael Stuart Brown, ``Nobel Lecture,''\emph{NobelPrize.org}. Winner-authored primary source; downloadable from the official Nobel Prize archive.\\
\url{https://www.nobelprize.org/prizes/medicine/1985/brown/lecture/}
\end{thebibliography}
